\section*{The exact price of endpoint exclusion}
\addcontentsline{toc}{section}{The exact price of endpoint exclusion}

The cavity Bayes floor can be maximal even in a two-bit example.  Let $X_i$ and $X_j$ be independent
uniform binary variables and define a desired edge label
\begin{equation}
Z_{ij}=\mathbf1[X_i=X_j].
\label{eq:book-endpoint-label}
\end{equation}
An endpoint-visible router predicts the label perfectly: it compares the two bits.  A strict cavity
router observes only
\begin{equation}
Y_{ij}=X_{[L]\setminus\{i,j\}},
\end{equation}
which is independent of $Z_{ij}$ in this construction.  Consequently,
\begin{equation}
\Pr(Z_{ij}=1\mid Y_{ij})=\frac12,
\qquad
e_{ij}^{\mathrm{cav}}=\frac12.
\label{eq:book-maximal-cavity-floor}
\end{equation}
No training procedure, teacher graph, model scale, or randomized decision rule can reduce the
strict inference error below one half without restoring endpoint information or adding a correlated
side channel.

The example exposes the semantic tradeoff cleanly.  The endpoint-visible router has perfect label
accuracy, but its support changes when either endpoint is intervened upon.  It is a valid dynamic
computational route and an invalid fixed conditional pair selector.  The cavity router preserves
the edge definition during endpoint substitution, but the desired label is not identifiable from
its allowed information.

Real proteins lie between these extremes.  Neighboring sequence context, predicted secondary
structure, templates, MSA-derived family evidence, or other modalities may make an edge partly
predictable without direct endpoint access.  The cavity posterior then moves away from one half and
the Bayes floor falls.  That gain should be attributed to the added observable evidence, not to an
architecture overcoming an information theorem.

This is why router performance and interaction semantics must be reported together.  A useful
comparison includes at least three variants:
\begin{equation}
\begin{aligned}
q_{ij}^{\mathrm{visible}}&=q(X_i,X_j,Y_{ij}),
&&\text{adaptive computational support},\\
q_{ij}^{\mathrm{cavity}}&=q(Y_{ij}),
&&\text{strict endpoint-invariant support},\\
q_{ij}^{\mathrm{fixed}}&=q(i,j),
&&\text{static candidate geometry}.
\end{aligned}
\label{eq:book-router-comparison}
\end{equation}
The first measures the attainable predictive frontier, the second the price of strict conditional
meaning, and the third the value of content-dependent routing itself.  No one variant is universally
correct; each licenses a different claim.
